\section{Related Work}

\subsection{In-Kernel Provenance and Information Flow}

ActPlane builds on a long line of provenance and information-flow systems.
PASS made provenance a storage-layer primitive~\cite{muniswamy2006pass}.
Hi-Fi and LPM used LSM hooks to capture trustworthy system
provenance~\cite{pohly2012hifi,bates2015lpm}. CamFlow captured whole-system
provenance streams over processes, files, and sockets~\cite{pasquier2017camflow}.

CamQuery is the closest prior system~\cite{pasquier2018camquery}. It executes
provenance queries inline in the kernel, propagates labels across process, file,
and network edges, and can prevent policy violations. ActPlane does not claim
that this mechanism is new. The difference is the domain and substrate:
ActPlane targets cooperative AI agents, compiles an agent-oriented source/sink
DSL to eBPF/BPF-LSM, and returns corrective feedback to the actor that violated
the rule.

TaintDroid, Dytan, libdft, and Panorama established dynamic taint tracking at
runtime, binary-instrumentation, or emulation layers~\cite{enck2010taintdroid,
clause2007dytan,kemerlis2012libdft,yin2007panorama}. ActPlane uses a much
coarser object-level taint. That coarseness is what allows kernel-wide coverage,
but it creates an over-taint risk that must be measured.

\subsection{eBPF and LSM Enforcement}

BPF-LSM, originally proposed through KRSI, lets eBPF programs attach to Linux
security hooks and return allow/deny verdicts~\cite{singh2019krsi,linuxbpflsm}.
ActPlane uses this as its pre-operation backend. Contemporary eBPF security
tools and proposals such as Tetragon, Tracee, and eBPF-PATROL demonstrate
practical event capture and enforcement, but generally evaluate per-event
predicates or lineage flags rather than typed cross-channel information
flow~\cite{ciliumtetragon,aquasecuritytracee,ghimire2025ebpfpatrol}.

\subsection{Agent Guardrails}

AgentSpec and Progent are the closest agent-policy systems at the tool
layer~\cite{wang2025agentspec,shi2025progent}. They provide deterministic
checks over tool names, arguments, or framework actions. ActPlane complements these systems by
moving the enforcement point below the tool API. If all effects pass through the
framework, tool-layer policies are sufficient and cheaper. If the agent can
execute arbitrary code, spawn subprocesses, or use direct system interfaces, the
OS layer is the robust boundary.

FIDES and CaMeL show that typed information-flow and capability separation are
useful abstractions for LLM agents~\cite{costa2025fides,debenedetti2025camel}. They operate
inside the agent loop or scaffolding. ActPlane applies the same style of
information-flow reasoning at the OS boundary.

\subsection{Agent Instruction Files}

Recent empirical work studies \texttt{CLAUDE.md}, \texttt{AGENTS.md}, and other
agent context files: their structure, content, maintenance, and effect on task
efficiency~\cite{chatlatanagulchai2025manifests,chatlatanagulchai2025agentreadmes,
santos2025claudeconfig,lulla2026agentsmd}. ActPlane's corpus study asks a
different question: which instructions are behavioral or flow invariants, and
which of those can be observed and enforced below the tool layer?
